%% =====================================================================
%%  SUPPLEMENTARY MATERIAL
%%
%%  Coarse Observability Across High-Order Finite-Volume
%%  Reconstructions: Universal Rank, Conditioning, and
%%  Dissipative Memory
%%
%%  Secondary material relocated from the main article's appendices
%%  (round-1 item P2, chain-independent length pass): the composed
%%  WENO5-LIN+SSP-RK3 coefficient table and the full one-at-a-time
%%  sensitivity analysis of the nonlinear ranks.
%%
%%  Cross-references to the main article are resolved with xr-hyper via
%%  \externaldocument{article}; build article.pdf first so that
%%  article.aux exists, then build this file.
%% =====================================================================
\documentclass[11pt,a4paper]{article}

\usepackage[a4paper,margin=1in]{geometry}
\usepackage{amsmath,amssymb,amsfonts}
\usepackage{mathtools}
\usepackage{booktabs}
\usepackage{array}
\usepackage{xcolor}
\usepackage{microtype}
\usepackage{xr-hyper}
\usepackage{hyperref}
\hypersetup{colorlinks=true,linkcolor=blue,citecolor=blue,urlcolor=blue,
  pdftitle={Supplementary Material: Coarse Observability Across High-Order Finite-Volume Reconstructions},
  pdfauthor={Antonis Polemitis, Ioannis Kokkinakis, Nicholas Christakis, Dimitris Drikakis}}
\usepackage[nameinlink]{cleveref}

% Label texts imported from article.aux use these article.tex macros.
\newcommand{\OL}{{\cal O}_L}
\newcommand{\Asch}{A^{[\rm sch]}_\lambda}
\DeclareMathOperator{\rank}{rank}

% Resolve \ref{...} to the main article's numbers.
\externaldocument{article}

% Supplementary numbering: S1, S2, ... for sections, tables, equations.
\renewcommand{\thesection}{S\arabic{section}}
\renewcommand{\thetable}{S\arabic{table}}
\renewcommand{\theequation}{S\arabic{equation}}

\title{Supplementary Material\\[2pt]
\large Coarse Observability Across High-Order Finite-Volume Reconstructions: Universal Rank, Conditioning, and Dissipative Memory}
\author{Antonis Polemitis \and Ioannis Kokkinakis \and Nicholas Christakis \and Dimitris Drikakis}
\date{}

\begin{document}
\maketitle

\noindent This Supplementary Material collects two blocks of secondary detail relocated from the main article to keep the article self-contained while shortening it: the exact coefficient polynomials of the composed WENO5-LIN + SSP-RK3 one-step operator, and the full one-at-a-time sensitivity analysis behind the nonlinear-rank robustness statement.
Both are machine-generated by \texttt{make\_tables.py} from the same exact rational-arithmetic representations and CSV artifacts as the in-article tables; section and equation numbers of the form ``6'', ``Table~6'', ``Appendix~A'' below refer to the main article.

\section{The SSP-RK3 composition}\label{app:rk3-coefficients}

For a linear autonomous update the Shu--Osher stages reduce to the classical third-order stability polynomial: with $A_{\rm FE}=I+Z$ and $Z=A_{\rm FE}-I$ the ($\lambda$-scaled) flux-difference part of Appendix~\ref{app:weno-coefficients} of the main article, \[ A_{\rm RK3} =I+Z+\tfrac12\,Z^{2}+\tfrac16\,Z^{3}, \] a circulant with support $s\in\{-6,\dots,9\}$, i.e.\ $(m_L,m_R)=(9,6)$ and total reach $w=15$.
\Cref{tab:coeff-rk3} lists the sixteen cubic coefficient polynomials.
The outermost coefficients $c_{9}=\lambda^{3}/162000$ and $c_{-6}=\lambda^{3}/48000$ are nonzero for every $\lambda>0$, so the reaches and collar counts of the pairing are $\lambda$-independent on $(0,1]$ (cf.\ Remark~\ref{rem:effective-support} of the main article).

\begin{table}[htbp]
\centering
\caption{Exact coefficient polynomials of the composed WENO5-LIN + SSP-RK3 one-step operator (sixteen cells, paired columns).}
\label{tab:coeff-rk3}
\begin{tabular}{@{}rl@{\qquad}rl@{}}
\toprule
$s$ & $c_s(\lambda)$                                                                       & $s$  & $c_s(\lambda)$                                                                            \\
\midrule
$9$ & $\tfrac{1}{162000}\lambda^{3}$                                                       & $1$  & $\lambda - \tfrac{31}{150}\lambda^{2} - \tfrac{49}{200}\lambda^{3}$                       \\
$8$ & $-\tfrac{1}{7200}\lambda^{3}$                                                        & $0$  & $1 - \tfrac{1}{3}\lambda - \tfrac{329}{720}\lambda^{2} + \tfrac{10211}{64800}\lambda^{3}$ \\
$7$ & $\tfrac{23}{14400}\lambda^{3}$                                                       & $-1$ & $-\tfrac{1}{2}\lambda + \tfrac{13}{60}\lambda^{2} + \tfrac{6253}{72000}\lambda^{3}$       \\
$6$ & $\tfrac{1}{1800}\lambda^{2} - \tfrac{961}{86400}\lambda^{3}$                         & $-2$ & $\tfrac{1}{20}\lambda + \tfrac{13}{120}\lambda^{2} - \tfrac{1849}{28800}\lambda^{3}$      \\
$5$ & $-\tfrac{1}{120}\lambda^{2} + \tfrac{121}{2400}\lambda^{3}$                          & $-3$ & $-\tfrac{1}{40}\lambda^{2} - \tfrac{9}{800}\lambda^{3}$                                   \\
$4$ & $\tfrac{31}{480}\lambda^{2} - \tfrac{427}{3000}\lambda^{3}$                          & $-4$ & $\tfrac{1}{800}\lambda^{2} + \tfrac{7}{1200}\lambda^{3}$                                  \\
$3$ & $\tfrac{1}{30}\lambda - \tfrac{47}{180}\lambda^{2} + \tfrac{9467}{43200}\lambda^{3}$ & $-5$ & $-\tfrac{1}{1600}\lambda^{3}$                                                             \\
$2$ & $-\tfrac{1}{4}\lambda + \tfrac{17}{30}\lambda^{2} - \tfrac{449}{9600}\lambda^{3}$    & $-6$ & $\tfrac{1}{48000}\lambda^{3}$                                                             \\
\bottomrule
\end{tabular}
\end{table}

\section{Sensitivity of the numerical ranks}\label{app:sensitivity}

The nonlinear ranks of Section~\ref{sec:exp-nonlinear} of the main article depend on declared numerical choices: the rank floor $\theta$, the ensemble size $M$, the random directions, the finite-difference step $\delta$ of the linearized ladder, and the physical amplitude $\varepsilon$.
Experiment E6s varies each knob one at a time around the E6 baseline ($\theta=10^{-3}$, $M=400$ plus the deep pair, seed $0$, $\delta=10^{-5}$, $\varepsilon=10^{-5}$) over all $21$ (scheme, regime) cells and compares, per axis setting, the leading-floor ladder $\rho^{\rm emp}_\theta(L)$ for $L\le5$ and the horizon-12 rank against the baseline; for the FD-step axis the compared quantity is the linearized ladder $\rho^{\rm lin}$.
\Cref{tab:sensitivity} counts the cells that change.

\begin{table}[htbp]
\centering
\caption{One-at-a-time sensitivity of the E6 ranks (E6s): number of the $21$ (scheme, regime) cells whose leading-floor ladder $\rho^{\rm emp}_\theta(L)$, $L\le5$, or horizon-12 rank differs from the axis baseline ($\theta=10^{-3}$, $M=400$, seed $0$, $\delta=10^{-5}$, $\varepsilon=10^{-5}$), and the largest change in the horizon-12 rank.
The FD-step rows compare the linearized ladders $\rho^{\rm lin}$ (tolerance $10^{-8}$) instead.
The $M$ axis uses nested direction prefixes of one bank; the seed axis redraws all $M=400$ directions.
Body machine-generated by \texttt{make\_tables.py} from \texttt{nonlinear\_sensitivity.csv}.}
\label{tab:sensitivity}
\begin{tabular}{@{}lcccc@{}}
\toprule
                              &           & \multicolumn{2}{c}{cells changed (of 21)} &                                     \\
\cmidrule(lr){3-4}
axis                          & setting   & ladder $L\le5$                            & $\rho(12)$ & $\max|\Delta\rho(12)|$ \\
\midrule
rank floor $\theta$           & $10^{-2}$ & 11                                        & 11         & 7                      \\
                              & $10^{-4}$ & 1                                         & 12         & 4                      \\
\addlinespace
ensemble size $M$             & 100       & 0                                         & 1          & 1                      \\
                              & 200       & 0                                         & 1          & 1                      \\
                              & 800       & 0                                         & 0          & 0                      \\
\addlinespace
seed                          & 1         & 0                                         & 1          & 1                      \\
                              & 2         & 0                                         & 0          & 0                      \\
\addlinespace
FD step $\delta$ (linearized) & $10^{-4}$ & 0                                         & 0          & 0                      \\
                              & $10^{-6}$ & 0                                         & 0          & 0                      \\
\addlinespace
amplitude $\varepsilon$       & $10^{-2}$ & 13                                        & 17         & 14                     \\
                              & $10^{-3}$ & 4                                         & 14         & 6                      \\
                              & $10^{-7}$ & 0                                         & 1          & 2                      \\
\bottomrule
\end{tabular}
\end{table}

The methodological knobs are inert.
Growing the ensemble from $100$ to $800$ members, redrawing all directions (seeds $1$, $2$), and moving the FD step across two orders of magnitude leave every reported ladder unchanged; the only motion anywhere is a single borderline cell (MUSCL superbee, smooth regime) whose horizon-12 rank flips between $6$ and $7$ at $M\le200$ and seed $1$ --- one direction sitting at the declared floor.
In particular, the plateau staircase of Table~\ref{tab:nonlinear} persists identically across every ensemble size and seed.
The rank floor behaves as a declared tolerance should: tightening it to $10^{-4}$ touches one ladder cell (the same superbee cell, $6\to7$ at $L=5$) but raises the horizon-12 rank in $12$ cells (by up to $4$) as near-floor directions get counted, and loosening it to $10^{-2}$ drops weakly exposed directions in $11$ cells --- which is why every rank in Section~\ref{sec:exp-nonlinear} of the main article is quoted at a stated floor rather than as an absolute.

The amplitude axis moves ranks only where the physics in Section~\ref{sec:exp-nonlinear} requires it.
At $\varepsilon=10^{-7}$, the ladders are identical to the baseline (one WENO-Z horizon-12 cell loses two near-floor directions).
At $\varepsilon=10^{-3}$ --- the predicted crossover $\varepsilon\approx\sqrt{\epsilon_w}$ for $\epsilon_w=10^{-6}$ --- exactly the WENO5-JS ladders (extremum and jump regimes) and the WENO5-M jump ladder leave the linear-budget shape, while WENO-Z, whose crossover lies far below, keeps its ladders.
At $\varepsilon=10^{-2}$, the smooth- and extremum-regime cells respond broadly, as in Table~\ref{tab:nonlinear}.
Throughout, the four MUSCL jump-regime (flat-plateau) ladders are unchanged at \emph{every} amplitude from $10^{-7}$ to $10^{-2}$: the scale-freeness of the positively homogeneous plateau response (Section~\ref{sec:empirical-rank} of the main article) is exact in the sweep.

\end{document}
